% БҮЛЭГ 5 — ЗАГВАР ХӨГЖҮҮЛЭЛТ БА PIPELINE

\chapter{Загвар хөгжүүлэлт ба дамжлага}
\label{Chapter5}

%-------------------------------------------------------------------------------
\section{MN-BERT сургалтын дамжлагын тойм}

Эцсийн нэг шатлалт MN-BERT болон харьцуулалтын хоёр шатлалт MN-BERT хувилбарыг
PyTorch 2.x болон Hugging Face Transformers 4.x сангийн орчинд сургасан. Нэг
шатлалт загвар нь энэ ажлын үндсэн шийдэл бөгөөд хоёр шатлалт хувилбар нь
хоёр шатлалт зохион байгуулалтын нөлөөг хэмжих харьцуулалтын хувилбар юм. Сургалтын
ерөнхий дамжлагыг Хүснэгт~\ref{tab:pipeline-flow}-д нэгтгэв.

\begin{table}[H]
\centering
\caption{MN-BERT сургалтын дамжлагын үе шатууд}
\label{tab:pipeline-flow}
\fittable{%
\begin{tabular}{@{}clp{8cm}@{}}
\toprule
\textbf{№} & \textbf{Алхам} & \textbf{Тайлбар} \\
\midrule
1 & Өгөгдөл унших          & CSV файлаас \texttt{pandas} ашиглан уншиж, сургалт, баталгаажуулалт, туршилтын хуваалтад ангилах \\
2 & Токенчлол              & \texttt{AutoTokenizer.from\_pretrained} дуудаж, \texttt{bert-base-mongolian-cased} загвараар токенчлох; эцсийн загварт \texttt{max\_length}=256 ашиглах \\
3 & Өгөгдлийн объект       & \texttt{input\_ids}, \texttt{attention\_mask}, \texttt{labels} тензоруудыг бэлтгэх \\
4 & Багц үүсгэх            & \texttt{DataLoader}-оор багц үүсгэх; сургалтын үед түүвэрлэлтийн тохиргоо ашиглаж, баталгаажуулалт ба туршилтад дарааллыг өөрчлөхгүй унших \\
5 & Урагш дамжуулалт       & \texttt{outputs = model(input\_ids, attention\_mask, labels)} хэлбэрээр гаралтыг тооцох \\
6 & Алдагдал тооцох        & Нэг шатлалт загварт фокус алдагдлын функц (Focal Loss) ($\gamma{=}2.0$) ба ангиллын жин; хоёр шатлалт хувилбарт жинтэй кросс-энтропи ашиглах \\
7 & Буцаан дамжуулалт      & \texttt{loss.backward()} ажиллуулж, градиентыг \texttt{max\_norm=1.0}-ээр хязгаарлах \\
8 & Жинг шинэчлэх          & AdamW оновчлогч, жингийн бууралтын тохиргоотойгоор параметр шинэчлэх \\
9 & Суралцах хурдын хуваарь & Эцсийн нэг шатлалт загварт косинус хуваарьтай халаалт (cosine warmup), хоёр шатлалт хувилбарт шугаман халаалттай хуваарь ашиглах \\
10 & Баталгаажуулалт       & Үе бүрийн дараа баталгаажуулалтын багцаар макро-F1 үзүүлэлтийг үнэлж, хамгийн сайн жинг хадгалах \\
11 & Шилдэг жинг хадгалах  & Шилдэг жинг хадгалж, эрт зогсоох нөхцөл ашиглах \\
\bottomrule
\end{tabular}}
\end{table}

Эцсийн нэг шатлалт загварт дөрвөн ангиллыг нэг загвараар шууд таамаглуулсан.
Харьцуулалтын хоёр шатлалт нэгдсэн дамжлагад эхлээд эерэг, саармаг, сөрөг
гэсэн ерөнхий шийдвэр гарч, сөрөг гэж шийдэгдсэн текстийг хортой сөрөг эсвэл
бүтээлч шүүмжлэл болгон нарийвчилсан. Таамаглал хийх үед энэ дамжлага нэг
эцсийн дөрвөн ангиллын гаралт үүсгэнэ.

%-------------------------------------------------------------------------------
\section{Оновчтой шийдэл: нэг шатлалт загварт шилжсэн нь}

Хоёр шатлалт туршилтаас шат дамжсан алдааны нөлөө илэрсэн: нэгдүгээр шатанд буруу
ангилагдсан текст хоёрдугаар шатанд зөв очих боломжгүй болдог. Иймээс дөрвөн
ангиллыг нэгэн зэрэг сурдаг \textbf{нэг шатлалт MN-BERT}-ийг эцсийн шийдэл
болгон сонгосон. Энэ хувилбар нь Focal Loss, ангиллын жин,
\texttt{WeightedRandomSampler} болон Cosine Warmup-ийн хослолоор хамгийн өндөр
гүйцэтгэл үзүүлсэн.

%-------------------------------------------------------------------------------
\section{Гиперпараметрийн тохиргоо ба оновчлол}

Эцсийн нэг шатлалт загварын гиперпараметрийн сонголтыг баталгаажуулалтын багц
дээрх макро-F1 үзүүлэлтээр оновчлов. Харьцуулалтын хоёр шатлалт нэгдсэн
дамжлагын тохиргоог нэг дамжлагын түвшинд товч нэгтгэв
(Хүснэгт~\ref{tab:hyperparams-mnbert}).

\begin{table}[H]
\centering
\caption{Сургалтын гиперпараметрийн тохиргоо}
\label{tab:hyperparams-mnbert}
\fittable{%
\begin{tabular}{@{}p{3.0cm}p{4.0cm}p{4.6cm}@{}}
\toprule
\textbf{Тохиргоо} & \textbf{Нэг шатлалт MN-BERT (эцсийн)} & \textbf{Хоёр шатлалт нэгдсэн дамжлага} \\
\midrule
Суралцах хурд              & $3\times10^{-5}$ & $3\times10^{-5}$ болон $2\times10^{-5}$ \\
Багцын хэмжээ              & 16   & 16   \\
Хамгийн сайн үе            & 8    & Дамжлагын бүрэлдэхүүн сургалтын логт хадгалсан үеүүд \\
Халаалтын харьцаа          & 0.10 & 0.06 \\
Жингийн бууралт            & 0.01 & 0.01 \\
Уналтын хувь (dropout)     & 0.15 & 0.15 \\
Оролтын дээд урт           & 256  & 256  \\
Алдагдлын функц            & Фокус алдагдлын функц (Focal Loss), $\gamma{=}2.0$ & Жинтэй кросс-энтропи (weighted cross-entropy) \\
Тэнцвэржүүлэлт             & $\alpha$ жин + тэнцвэржүүлсэн түүвэрлэлт & Ангиллын жин + тэнцвэржүүлсэн түүвэрлэлт \\
Оновчлогч                  & AdamW & AdamW \\
Суралцах хурдын хуваарь    & Косинус хуваарьтай халаалт (cosine warmup) & Шугаман халаалттай хуваарь \\
Эрт зогсоох нөхцөл         & 3 & 5 \\
Санамсаргүй эх утга        & 42 & 42 \\
\bottomrule
\end{tabular}}
\end{table}

Суралцах хурдыг торон хайлтаар (grid search) шалгахад $3\times10^{-5}$,
багцын хэмжээ 16 байх нь баталгаажуулалтын багц дээр хамгийн сайн макро-F1
утгыг үзүүлсэн.

%-------------------------------------------------------------------------------
\section{Сургалтын дэд бүтэц ба орчин}

Загваруудыг сургах болон турших ажлыг локал тооцооллын орчинд гүйцэтгэсэн.
Техник хангамжийн хувьд \textbf{AMD Radeon RX 9070 GRE 12GB} график картыг
ашигласан бөгөөд \texttt{torch-directml} санг ашиглан DirectX-ээр дамжуулан
GPU тооцооллыг идэвхжүүлсэн.

%-------------------------------------------------------------------------------
\section{Ангиллын тэнцвэржүүлэлт}

Бүлэг~\ref{Chapter4}-д үзүүлсэн тэнцвэргүй байдал (эерэг 6.55\% хамгийн цөөн)
нь цөөнх ангиллын F1-ийг доройтуулж болзошгүй тул загвар тус бүрд тохирсон
тэнцвэржүүлэлт хийв.

\textbf{Суурь загваруудад} SMOTE ($k\!=\!5$) ашиглан цөөнх ангиллын өгөгдлийг
нийлэгээр нэмэгдүүлэв.

\textbf{Нэг шатлалт MN-BERT-д} фокус алдагдлын функц (Focal Loss) ($\gamma{=}2.0$)-ийг урвуу
давтамжийн жин $\alpha_c = N/(K\cdot n_c)$ (эерэг $3.707$, саармаг $0.835$,
бүтээлч $0.570$, хортой $1.286$) болон \texttt{WeightedRandomSampler}-тэй
хослуулав. Үр нөлөөг бүлэг~\ref{Chapter6}-ийн макро-F1 болон
Хүснэгт~\ref{tab:dq-progression}-аар үнэлэв.

%-------------------------------------------------------------------------------
\section{Gradio интерфейс ба дамжлагын нэгтгэл}


Эцсийн системийн интерфейсийг Gradio сан ашиглан гурван хэсэгтэйгээр
хөгжүүлсэн: \textit{Ангилах} (нэг шатлалт MN-BERT $83.26\%$ ба хоёр шатлалт
нэгдсэн дамжлага $78.94\%$ сонголттой, магадлалын баганан дүрслэлтэй),
\textit{Харьцуулах} (хоёр
архитектурыг зэрэг ажиллуулж шат дамжсан алдааны нөлөөг харах), \textit{Үнэлгээ}
(төөрөлдлийн матриц, үнэлгээний үзүүлэлтүүд). Загваруудыг шаардлагатай үед ачаалах
горимоор AMD DirectML дээр таамаглал хийлгэсэн.

Мөн загварыг гаднын системээс туршилтын байдлаар дуудах боломжийг шалгахын
тулд FastAPI-д суурилсан API загварчилсан хувилбар боловсруулав.
\texttt{/predict} төгсгөл нь нэг сэтгэгдлийн ангилал, \texttt{/batch} нь олон
мөрийн багц таамаглал, \texttt{/audit} нь аудитын маягийн нэгтгэсэн гаралтыг
буцаана. API нь загваруудыг шаардлагатай үед ачаалах горим (lazy loading) болон
архитектур сонгох параметрээр нэг шатлалт, хоёр шатлалт хувилбаруудыг туршилтын
зорилгоор дэмжинэ. Энэ хэсэг нь үйлдвэрлэлийн бүрэн систем бус, судалгааны үр
дүнг шалгах API прототип юм.

Таамаглалын хурдыг хэмжихийн тулд 5 төрлийн бодит сэтгэгдэл дээр $N{=}20$
удаагийн давталттай жишиг туршилт явуулав
(Хүснэгт~\ref{tab:latency}).

\begin{table}[H]
\centering
\caption{Таамаглалын хурдны хэмжилт}
\label{tab:latency}
\fittable{%
\begin{tabular}{@{}lccc@{}}
\toprule
\textbf{Хувилбар} & \textbf{Дундаж (мс)} & \textbf{P95 (мс)} & \textbf{Хоёрдугаар шат идэвхжсэн} \\
\midrule
Нэг шатлалт   & $\mathbf{8.7}$  & $8.5$  & --- \\
Хоёр шатлалт дамжлага & $12.6$ & $16.2$ & $60\%$ \\
\midrule
\textbf{Хурдны зөрүү} & \multicolumn{3}{l}{Нэг шатлалт загвар хоёр шатлалт нэгдсэн дамжлагаас $\mathbf{30.8\%}$ хурдан} \\
\bottomrule
\end{tabular}}
\end{table}

Нэг шатлалт загвар үргэлж нэг урагш дамжуулалт хийдэг бол хоёр шатлалт нэгдсэн дамжлага
оролтын $60\%$-д хоёрдугаар шатыг нэмж ажиллуулдаг тул дунджаар $30.8\%$ удаан
байна. Нэг шатлалт загварын дундаж $8.7$~мс нь бодит цагийн ${\leq}100$~мс
шаардлагыг хангаж байна (ХО-5: ${\geq}20\%$ бууралт --- 5 оноо; ХО-19:
${\leq}100$~мс --- биелсэн).

\section{Давтагдах чадвар ба туршилтын бүртгэл}

Туршилтын үр дүнг давтах боломжийг хангах үүднээс санамсаргүй тооны үүсгүүрийн
эх утгыг 42 болгон тогтворжуулж, Git/GitHub ашиглан хувилбарын хяналт хийв. Мөн
\texttt{requirements.txt} файл, сургалтын явцын лог (loss/F1 JSON),
төөрөлдлийн матриц болон хамгийн сайн F1 утгатай жингийн файлуудыг хадгалсан.
Энэ нь Ш-6 шаардлагыг бүрэн хангана.
%-------------------------------------------------------------------------------
\section{Бүлгийн дүгнэлт}

Энэ бүлэгт MN-BERT сургалтын дамжлага, гиперпараметр, тэнцвэржүүлэлтийн арга
(SMOTE, фокус алдагдлын функц (Focal Loss), ангиллын жин,
\texttt{WeightedRandomSampler}), Gradio интерфейс, API прототип болон туршилтыг
дахин гүйцэтгэх нөхцөлийг баримтжуулсан. Дараагийн бүлэгт туршилтын үр дүнг
авч үзнэ.
